%=============================================================================!
% 2.2  FRAMES OF REFERENCE                                                    !
%=============================================================================!
\section{Frames of reference}
\label{sec:ne-frames}

Every position in Chapter~1 was measured from some origin, along some
axes, with times read from some clock. That choice is called a frame of
reference. \Cref{sec:mo-atoms} used the word frame for a saved snapshot
of positions, so here we always say frame of reference in full. Two
observers moving relative to each other use different frames of
reference and describe the same motion differently. This section works
out how their descriptions are related, which tells us from which frames
of reference the first law can hold.

\subsection*{One throw seen from two frames of reference}

A train runs along a straight, level track at a constant velocity
$\vb{u}$. We take the $x$ axis along the track in the direction of travel
and the $y$ axis vertically upwards, so that $\vb{u} = (5,
0)$\,\si{\metre\per\second}. A passenger throws a ball straight up, as
the passenger sees it, at \SI{6}{\metre\per\second}. To the passenger the
ball rises and falls along a vertical line and lands back in the hand
(\cref{fig:ne-frames}a). Someone standing on the platform sees the same
ball travel forwards with the train while it rises and falls, and so
trace out a parabola (\cref{fig:ne-frames}b). Both descriptions are
correct; they differ because they are made from different frames of
reference.

\begin{figure}[tb]
\centering
\includegraphics{ch02_newton/frames}
\caption{A ball thrown straight up at \SI{6}{\metre\per\second} by a
passenger in a train moving at $\vb{u} = (5, 0)$\,\si{\metre\per\second},
at five equally spaced times from the throw to the catch
\SI{1.224}{\second} later. (a) In the frame of reference of the train
the ball moves on a vertical line; it passes the same heights on the way
up and down, so the five positions show as three. Arrows give its
acceleration. (b) In the frame of reference of the platform it follows a
parabola, with velocity (teal) and acceleration (dark red); squares mark
the passenger's hand. (c) The velocities at the second instant, $t = 1.224/4 =
\SI{0.306}{\second}$, when the passenger sees the ball rising at $6 -
9.80665\times0.306 = \SI{3}{\metre\per\second}$: the platform sees $\vb{v}$, which is the velocity
$\vb{v}' = (0, 3)$\,\si{\metre\per\second} seen in the train plus the
velocity $\vb{u}$ of the train.}
\label{fig:ne-frames}
\end{figure}

To relate the two descriptions, the platform observer measures positions
from a fixed mark on the platform, the passenger measures them from the
hand, and both sets of axes point the same way. The hand moves at the
constant velocity $\vb{u}$, so by \cref{eq:mo-constant} with zero
acceleration it is at $\vb{r}_0 + \vb{u}t$ at time $t$, where $\vb{r}_0$
is its position at $t = 0$. To go from the mark to the ball we can go
first from the mark to the hand and then from the hand to the ball
(\cref{fig:ne-positions}), and adding these two vectors gives the
position $\vb{r}$ of the ball seen from the platform in terms of its
position $\vb{r}'$ seen from the hand,
\begin{equation*}
   \vb{r} = \vb{r}_0 + \vb{u}t + \vb{r}' .
\end{equation*}
Subtracting $\vb{r}_0 + \vb{u}t$ from both sides gives the position seen
in the train,
\begin{equation}
   \vb{r}' = \vb{r} - \vb{r}_0 - \vb{u}t .
   \label{eq:ne-galilean}
\end{equation}
Here we have assumed that both observers read the same time $t$ from
their clocks, which is true to an excellent approximation for every speed
far below that of light, about \SI{3e8}{\metre\per\second}, and so for
everything in this book.

\begin{figure}[htb]
\centering
\begin{tikzpicture}[scale=1.15, line cap=round,
   vec/.style={-{Stealth[length=5pt]}, line width=0.9pt}]
   \coordinate (O) at (0,0);
   \coordinate (H) at (3.2,0.35);
   \coordinate (B) at (3.9,2.0);
   \draw[vec] (O) -- (H) node[midway, below] {\small $\vb{r}_0 + \vb{u}t$};
   \draw[vec] (H) -- (B) node[midway, right] {\small $\vb{r}'$};
   \draw[vec, accent] (O) -- (B) node[midway, above left] {\small $\vb{r}$};
   \fill (O) circle (1.6pt) node[below left] {\small mark on the platform};
   \fill (H) circle (1.6pt) node[below right] {\small hand};
   \fill[accent] (B) circle (2pt) node[above right] {\small ball};
\end{tikzpicture}
\caption{The position $\vb{r}$ of the ball seen from the platform is the
position $\vb{r}_0 + \vb{u}t$ of the hand plus the position $\vb{r}'$ of
the ball seen from the hand.}
\label{fig:ne-positions}
\end{figure}

\subsection*{Velocities and accelerations}

Differentiating \cref{eq:ne-galilean} with respect to time relates the
velocities. The starting point $\vb{r}_0$ is fixed, so its derivative is
zero, and the derivative of $\vb{u}t$ is $\vb{u}$ because $\vb{u}$ is
constant:
\begin{align*}
   \frac{\dd\vb{r}'}{\dd t}
   &= \frac{\dd\vb{r}}{\dd t} - \frac{\dd\vb{r}_0}{\dd t}
      - \frac{\dd(\vb{u}t)}{\dd t}
      && \text{differentiating each term}\\
   &= \vb{v} - \vb{0} - \vb{u}
      && \text{$\vb{r}_0$ fixed, $\vb{u}$ constant,}
\end{align*}
so that
\begin{equation}
   \vb{v}' = \vb{v} - \vb{u} .
   \label{eq:ne-velocity-frames}
\end{equation}
The velocity seen in the train is the velocity seen from the platform,
less the velocity of the train. Differentiating once more, $\vb{u}$ is
constant and drops out entirely, and the two observers find the same
acceleration,
\begin{equation*}
   \vb{a}' = \vb{a} .
\end{equation*}

For the throw, the passenger sees the ball leave the hand with velocity
$\vb{v}' = (0, 6)$\,\si{\metre\per\second}, so the platform sees it leave
with $\vb{v} = \vb{v}' + \vb{u} = (5, 6)$\,\si{\metre\per\second}, a
speed of $(5^2 + 6^2)^{1/2} = \SI{7.810}{\metre\per\second}$.
\Cref{fig:ne-frames}(c) draws this sum at a later instant. Both observers
see the same acceleration, $(0, -g)$, and, because $\vb{u}$ has no
vertical component, the same vertical starting velocity,
\SI{6}{\metre\per\second}, so they see the same vertical motion. By
\cref{eq:mo-constant,eq:mo-v2}, the ball rises for $6/g =
\SI{0.612}{\second}$ to a greatest height of $6^2/2g = \SI{1.835}{\metre}$
and falls back in the same time, landing after \SI{1.224}{\second}. In
that time the platform observer sees it travel $5\times1.224 =
\SI{6.12}{\metre}$ along the track, which is exactly how far the hand has
moved, so the passenger catches it.

\subsection*{Inertial frames of reference}

Positions and velocities depend on the frame of reference, but
accelerations do not, for any two frames of reference moving at constant
velocity relative to each other. The agreement fails when one of them
accelerates. If the train brakes, the hand no longer moves at constant
velocity. We write its position as some function $\vb{s}(t)$, which for
the steady train was $\vb{r}_0 + \vb{u}t$. The same subtraction then
gives $\vb{r}' = \vb{r} - \vb{s}(t)$, and differentiating twice,
\begin{equation*}
   \vb{a}' = \vb{a} - \frac{\dd^2\vb{s}}{\dd t^2} .
\end{equation*}
The two observers now disagree about accelerations, by exactly the
acceleration of the train. Picture a cup on a smooth table in the train.
If the train slows at \SI{2}{\metre\per\second\squared}, its acceleration
is $\dd^2\vb{s}/\dd t^2 = (-2, 0)$\,\si{\metre\per\second\squared},
pointing backwards. Nothing pushes the cup along the track, so from the
platform $\vb{a} = \vb{0}$, and the passenger sees $\vb{a}' = \vb{0} -
(-2, 0) = (2, 0)$\,\si{\metre\per\second\squared}: the cup slides
forwards with nothing pushing it, which breaks the first law, while from the platform it simply
keeps moving at its old velocity as the train slows beneath it, as the
first law says.

The first law therefore cannot hold in every frame of reference, and we
turn this round into a definition. A frame of reference in which a body
with no net force on it moves at constant velocity is called inertial,
from inertia, the tendency of a body to keep the velocity it has. What
experiment adds is that such frames of reference exist. A frame of
reference fixed to the ground is very nearly inertial. The Earth turns
once in \SI{86164}{\second} measured against the distant stars. This is
a little less than the \SI{86400}{\second} of a 24-hour day, which is
measured against the Sun: the Earth's yearly journey round the Sun adds
one turn against the stars each year, so $86400\times365.24/366.24 =
\SI{86164}{\second}$. Its angular speed is therefore $\omega = 2\pi/86164
= \SI{7.29e-5}{\radian\per\second}$. In a frame of reference with its
origin at the Earth's centre and its axes pointing at fixed distant
stars, a point on the equator, carried round a circle of radius $R =
\SI{6.38e6}{\metre}$, has by \cref{eq:mo-circle} an acceleration of only
$\omega^2 R = (7.29\times10^{-5})^2\times6.38\times10^6 =
\SI{0.034}{\metre\per\second\squared}$, about \SI{0.35}{\percent} of
$g$. By $\vb{a}' = \vb{a} - \dd^2\vb{s}/\dd t^2$, a body with no net
force on it therefore seems, from the ground, to accelerate by at most
about this much; the turning of the ground's axes adds effects that are
smaller still at everyday speeds.
Since accelerations agree between frames of reference moving at constant
velocity relative to each other, every frame of reference moving at
constant velocity relative to the ground is very nearly inertial as well. All three
of Newton's laws are stated for inertial frames of reference, and from
here on every motion is described in one.

\notebook{02}{change the train's speed and watch both views of the throw}

\begin{selfcheck}
A passenger walks towards the back of the train at \SI{1}{\metre\per\second}
while the train moves forwards at \SI{5}{\metre\per\second}. What is the
passenger's velocity seen from the platform? Is a spinning roundabout an
inertial frame of reference?
\end{selfcheck}
